% ============================================================
% content/solutions.tex —— 2023 湖南卷物理 · 参考答案与详细解析
% 解析类型：source_embedded（转录自来源站 Word 版内嵌解析，独立复核）
% 复核修正见 metadata.yaml 的 notes
% ============================================================

\begin{answer}
\gksolhead{1}{单项选择题}
\gkanswer{答案 A}
A．相同质量的核燃料，轻核聚变的质量亏损比例更大，释放的核能更多，A 正确。

B．氘氚聚变放出的应是中子，方程为 ${}_{1}^{2}\mathrm{H}+{}_{1}^{3}\mathrm{H}\to{}_{2}^{4}\mathrm{He}+{}_{0}^{1}\mathrm{n}$。题给方程左边电荷数为 $1+1=2$，右边为 $2-1=1$，电荷数不守恒，B 错误。

C．铀 235 是重核裂变的燃料，聚变燃料是氘、氚等轻核，C 错误。

D．聚变放出能量，必然伴随质量亏损，D 错误。
\end{answer}

\begin{answer}
\gksolhead{2}{单项选择题}
\gkanswer{答案 B}
两谷粒在空中均只受重力，加速度都等于重力加速度 $g$，A 错误。

设 $O$、$P$ 的高度差为 $h$。谷粒 1 做平抛运动，有
\[
  h=\frac{1}{2}gt_{1}^{2}.
\]
谷粒 2 做斜上抛运动，取竖直向下为正方向，有
\[
  h=-v_{2y}t_{2}+\frac{1}{2}gt_{2}^{2},
\]
其中 $v_{2y}>0$ 为初速度的竖直分量。两式相减得
\[
  \frac{1}{2}g\left(t_{2}^{2}-t_{1}^{2}\right)=v_{2y}t_{2}>0,
\]
故 $t_{2}>t_{1}$，C 错误。

两谷粒从 $O$ 到 $P$ 的位移相同，而 $t_{2}>t_{1}$，由平均速度的定义可知谷粒 2 的平均速度较小，D 错误。

水平方向位移相同，即 $v_{1}t_{1}=v_{2x}t_{2}$。由 $t_{2}>t_{1}$ 得 $v_{2x}<v_{1}$。谷粒 2 在最高点的速度就是其水平分速度 $v_{2x}$，故小于 $v_{1}$，B 正确。
\end{answer}

\begin{answer}
\gksolhead{3}{单项选择题}
\gkanswer{答案 C}
由题图乙可知，波源振动的周期 $T=\SI{4}{\second}$，振幅 $A=\SI{2}{\centi\metre}$；波长 $\lambda=\SI{4}{\metre}$，故波速
\[
  v=\frac{\lambda}{T}=\frac{4}{4}\,\si{\metre\per\second}=\SI{1}{\metre\per\second},
\]
A 错误。

$C$ 到 $D$ 的距离为 \SI{3}{\metre}，$A$、$B$ 到 $D$ 的距离均为
\[
  \sqrt{4^{2}+3^{2}}\,\si{\metre}=\SI{5}{\metre}.
\]
波传到 $D$ 的最短时间为 \SI{3}{\second}，故 $t=\SI{2}{\second}$ 时 $D$ 处质点尚未开始振动，B 错误。

$t=\SI{4.5}{\second}$ 时，$A$、$B$ 的波尚未传到 $D$（需 \SI{5}{\second}），只有 $C$ 的波使 $D$ 振动了 \SI{1.5}{\second}，即 $\tfrac{3}{8}T$。由题图乙，$t=\SI{1.5}{\second}$ 时波源正从正向最大位移处向平衡位置运动，速度沿 $-y$ 方向；$D$ 处质点重复波源的振动，故也向 $y$ 轴负方向运动，C 正确。

$t=\SI{6}{\second}$ 时，$C$ 的波使 $D$ 振动了 \SI{3}{\second}，即 $\tfrac{3}{4}T$，位移为 \SI{-2}{\centi\metre}；$A$、$B$ 的波各使 $D$ 振动了 \SI{1}{\second}，即 $\tfrac{1}{4}T$，位移各为 \SI{+2}{\centi\metre}。由波的叠加原理，$D$ 处质点的位移
\[
  y=-2+2+2=+\SI{2}{\centi\metre},
\]
与平衡位置相距 \SI{2}{\centi\metre}，不是 \SI{6}{\centi\metre}，D 错误。
\end{answer}

\begin{answer}
\gksolhead{4}{单项选择题}
\gkanswer{答案 B}
恒星自转，表面各点随恒星自转做圆周运动，所需的向心力不同，故同一恒星表面各处的重力加速度不完全相同，A 错误。

在恒星表面两极处，万有引力等于重力：
\[
  G\frac{Mm}{R^{2}}=mg,\qquad g=\frac{GM}{R^{2}}.
\]
坍缩后质量 $M$ 不变、半径 $R$ 减小，故 $g$ 增大，B 正确。

第一宇宙速度
\[
  v=\sqrt{\frac{GM}{R}},
\]
$R$ 减小则 $v$ 增大，C 错误。

逃逸速度 $v_{\text{逃}}=\sqrt{2}\,v=\sqrt{\dfrac{2GM}{R}}$。由 $M=\rho\cdot\dfrac{4}{3}\pi R^{3}$ 得
\[
  R=\left(\frac{3M}{4\pi\rho}\right)^{1/3},
\]
代入得
\[
  v_{\text{逃}}=\sqrt{2GM}\left(\frac{4\pi\rho}{3M}\right)^{1/6}
  \propto M^{1/3}\rho^{1/6}.
\]
质量相同时，密度越大逃逸速度越大；中子星的密度大于白矮星，故中子星的逃逸速度更大，D 错误。
\end{answer}

\begin{answer}
\gksolhead{5}{单项选择题}
\gkanswer{答案 D}
设 $PQ_{1}=h$。由几何关系，三个点电荷到 $P$ 的距离分别为
\[
  r_{1}=h,\qquad r_{2}=\frac{h}{\sin60^{\circ}}=\frac{2h}{\sqrt{3}},\qquad
  r_{3}=\frac{h}{\sin30^{\circ}}=2h.
\]
以 $Q_{1}\to Q_{3}$ 为 $x$ 轴正方向、$Q_{1}\to P$ 为 $y$ 轴正方向，三个点电荷在 $P$ 点产生的场强分别沿 $PQ_{1}$、$Q_{2}P$、$Q_{3}P$ 方向，与 $x$ 轴正方向的夹角依次为 $90^{\circ}$、$120^{\circ}$、$150^{\circ}$。$P$ 点合场强为零，故
\[
  E_{x}=k\left[\frac{Q_{2}}{r_{2}^{2}}\cos120^{\circ}
            +\frac{Q_{3}}{r_{3}^{2}}\cos150^{\circ}\right]=0,
\]
\[
  E_{y}=k\left[\frac{Q_{1}}{r_{1}^{2}}
            +\frac{Q_{2}}{r_{2}^{2}}\sin120^{\circ}
            +\frac{Q_{3}}{r_{3}^{2}}\sin150^{\circ}\right]=0.
\]
代入 $r_{1}$、$r_{2}$、$r_{3}$，得
\[
  -\frac{3}{8}Q_{2}-\frac{\sqrt{3}}{8}Q_{3}=0
  \;\Longrightarrow\;
  3Q_{2}+\sqrt{3}Q_{3}=0,
\]
\[
  Q_{1}+\frac{3\sqrt{3}}{8}Q_{2}+\frac{1}{8}Q_{3}=0.
\]
解得
\[
  Q_{2}=-\frac{4\sqrt{3}}{3}Q_{1},\qquad Q_{3}=4Q_{1}.
\]
取 $Q_{1}=q$（$q>0$），得 $Q_{2}=-\dfrac{4\sqrt{3}}{3}q$，$Q_{3}=4q$，D 正确。

A、B 中三个电荷电性相同，三个场强不可能相互抵消；C 中 $Q_{1}=-q$、$Q_{2}=\sqrt{2}q$、$Q_{3}=-q$ 代入 $x$ 分量式得 $3\sqrt{2}q-\sqrt{3}q\ne0$，均错误。
\end{answer}

\begin{answer}
\gksolhead{6}{单项选择题}
\gkanswer{答案 D}
在区域Ⅰ中，粒子沿直线 $AC$ 运动，电场力与洛伦兹力平衡：
\[
  qE=qvB_{1},\qquad v=\frac{E}{B_{1}}.
\]
进入区域Ⅱ后，粒子在洛伦兹力作用下做匀速圆周运动。由左手定则，正电荷、$B_{2}$ 垂直纸面向外、速度沿 $AC$ 方向时，洛伦兹力指向下方，故圆心 $O'$ 在过 $C$ 的竖直线下方。

设 $CF=a$，则 $GF$ 到 $C$ 的水平距离为 $\dfrac{\sqrt{2}}{2}a$。粒子从 $CF$ 中点 $M$ 射出时，$CM=\dfrac{a}{2}$，以 $C$ 为原点、$AC$ 为 $x$ 轴，则 $M$ 的坐标为 $\left(\dfrac{\sqrt{2}}{4}a,\,-\dfrac{\sqrt{2}}{4}a\right)$，圆心 $O'$ 的坐标为 $(0,\,-r)$。由 $|O'M|=r$ 得
\[
  \left(\frac{\sqrt{2}}{4}a\right)^{2}
  +\left(r-\frac{\sqrt{2}}{4}a\right)^{2}=r^{2},
\]
解得
\[
  r=\frac{\sqrt{2}}{4}a.
\]
轨迹对应的圆心角为 $90^{\circ}$，故
\[
  t_{0}=\frac{T}{4}=\frac{\pi m}{2qB_{2}}.
\]
五种改变方式下粒子在区域Ⅱ中的轨迹如图甲—戊所示。

A．$B_{1}$ 变为 $2B_{1}$ 时，$v_{1}=\dfrac{E}{2B_{1}}=\dfrac{v}{2}$，$r_{1}=\dfrac{r}{2}$，粒子仍从 $CF$ 边射出，圆心角仍为 $90^{\circ}$，$t=t_{0}$，A 错误。

B．$E$ 变为 $2E$ 时，$v_{2}=2v$，$r_{2}=2r=\dfrac{\sqrt{2}}{2}a$，粒子恰好从 $F$ 点射出，圆心角仍为 $90^{\circ}$，$t=t_{0}$，B 错误。

C．$B_{2}$ 变为 $\dfrac{\sqrt{3}}{4}B_{2}$ 时，
\[
  r_{3}=\frac{mv}{\frac{\sqrt{3}}{4}qB_{2}}=\frac{4}{\sqrt{3}}r=\sqrt{\frac{2}{3}}\,a.
\]
粒子从 $GF$ 边射出。设圆心角为 $\theta$，圆心到 $GF$ 的距离为 $\dfrac{\sqrt{2}}{2}a$，则
\[
  \sin\theta=\frac{\frac{\sqrt{2}}{2}a}{r_{3}}=\frac{\sqrt{3}}{2},
  \qquad \theta=60^{\circ}.
\]
故
\[
  t_{3}=\frac{1}{6}\cdot\frac{2\pi m}{q\cdot\frac{\sqrt{3}}{4}B_{2}}
      =\frac{8\sqrt{3}}{9}\,t_{0}\ne\frac{t_{0}}{2},
\]
C 错误。

D．$B_{2}$ 变为 $\dfrac{\sqrt{2}}{4}B_{2}$ 时，
\[
  r_{4}=\frac{mv}{\frac{\sqrt{2}}{4}qB_{2}}=\frac{4}{\sqrt{2}}r=a,
\]
粒子从 $GF$ 边射出，且
\[
  \sin\theta'=\frac{\frac{\sqrt{2}}{2}a}{a}=\frac{\sqrt{2}}{2},
  \qquad \theta'=45^{\circ},
\]
故
\[
  t_{4}=\frac{1}{8}\cdot\frac{2\pi m}{q\cdot\frac{\sqrt{2}}{4}B_{2}}
      =\sqrt{2}\,t_{0},
\]
D 正确。

\gkfig{figures/s6-cases}
\end{answer}

\begin{answer}
\gksolhead{7}{多项选择题}
\gkanswer{答案 BC}
光从水中射向空气，入射角 $i=90^{\circ}-\alpha$。发生全反射的临界角 $C$ 满足
\[
  n=\frac{1}{\sin C}.
\]
当 $\alpha=41^{\circ}$ 时恰好发生全反射，故
\[
  C=90^{\circ}-41^{\circ}=49^{\circ},\qquad n=\frac{1}{\sin49^{\circ}},
\]
A 错误、B 正确。

当 $\alpha=60^{\circ}$ 时，$i=30^{\circ}<C$，光能射出水面。由折射定律
\[
  n\sin i=\sin r,
\]
且 $n>1$，得 $\sin r=n\sin30^{\circ}>\sin30^{\circ}$，故 $r>30^{\circ}$。出射光与法线夹角为 $r$，与水面夹角为 $90^{\circ}-r<60^{\circ}$，C 正确、D 错误。

\gkfig{figures/s7-rays}
\end{answer}

\begin{answer}
\gksolhead{8}{多项选择题}
\gkanswer{答案 AD}
小球恰好到达 $C$ 点，说明 $v_{C}=0$。$A$、$C$ 的高度差为 $2R$，由机械能守恒定律
\[
  \frac{1}{2}mv_{0}^{2}=mg\cdot2R,
\]
解得
\[
  v_{0}=\sqrt{4gR}=2\sqrt{gR},
\]
C 错误。

$A$ 到 $C$ 的过程中，小球竖直分速度由 $v_{0}\sin\theta$ 逐渐减小到 $0$，重力的功率 $P=mgv_{y}$ 逐渐减小，不是保持不变，B 错误。

设小球与 $O$ 的连线与竖直方向夹角为 $\alpha$。在圆弧上，沿半径方向有
\[
  mg\cos\alpha-F_{N}=m\frac{v^{2}}{R}.
\]
从 $B$ 到 $C$，$\alpha$ 减小、$\cos\alpha$ 增大，同时 $v$ 减小，故 $F_{N}$ 增大；由牛顿第三定律，小球对轨道的压力逐渐增大，A 正确。

在 $B$ 点，若 $v_{B}^{2}>gR\cos\theta$，则所需向心力大于重力沿半径方向的分力，轨道对小球的弹力将减为零，小球从 $B$ 点脱离轨道。$v_{0}$ 增大时 $v_{B}$ 增大，故小球有可能从 $B$ 点脱离轨道，D 正确。
\end{answer}

\begin{answer}
\gksolhead{9}{多项选择题}
\gkanswer{答案 AC}
皮带不打滑，两轮边缘的线速度大小相等：
\[
  \omega r_{\text{大}}=\omega' r_{\text{小}},\qquad \omega'=4\omega,
\]
即线圈转动的角速度为 $4\omega$，A 正确。

线圈电动势的最大值
\[
  E_{m}=nBL^{2}\cdot4\omega=4nBL^{2}\omega,
\]
有效值
\[
  E=\frac{E_{m}}{\sqrt{2}}=2\sqrt{2}nBL^{2}\omega.
\]
线圈与灯泡的阻值均为 $R$，灯泡两端电压
\[
  U=\frac{R}{R+R}E=\sqrt{2}nBL^{2}\omega,
\]
B 错误。

漆包线总长变为原来的 2 倍时，电阻变为 $2R$，匝数变为 $2n$。此时
\[
  E_{m}'=2nBL^{2}\cdot4\omega,\qquad E'=\frac{E_{m}'}{\sqrt{2}}=4\sqrt{2}nBL^{2}\omega,
\]
灯泡两端电压
\[
  U'=\frac{R}{R+2R}E'=\frac{4\sqrt{2}}{3}nBL^{2}\omega,
\]
C 正确。

小轮半径变为原来的 2 倍时，小轮（及线圈）的角速度变为
\[
  \omega''=\frac{\omega r}{2r}=2\omega,
\]
电动势减半，灯泡变暗，D 错误。
\end{answer}

\begin{answer}
\gksolhead{10}{多项选择题}
\gkanswer{答案 CD}
设杆对球的弹力大小为 $F_{T}$。$A$ 球在竖直方向只受重力和杆的作用力，故
\[
  F_{T}\cos\theta=mg,
\]
杆处于压缩状态：杆对 $A$ 的作用力沿杆向左上，对 $B$ 的作用力沿杆向右下。$B$ 球对底面的压力
\[
  N_{B}=mg+F_{T}\cos\theta=2mg,
\]
故 $B$ 球所受最大静摩擦力为 $2\mu mg$。

A．$B$ 球不受摩擦力时，其水平方向只受杆的分力：
\[
  F_{T}\sin\theta=ma,\qquad a=g\tan\theta.
\]
对整体 $F=4ma=4mg\tan\theta$，A 错误。

B、C．推力向左时，整体加速度 $a$ 向左。对 $A$，水平方向
\[
  F_{T}\sin\theta-N_{A}=ma,\qquad N_{A}=m\left(g\tan\theta-a\right)\ge0,
\]
故 $a\le g\tan\theta$。对 $B$，摩擦力向左，有
\[
  f-mg\tan\theta=ma,\qquad f=m\left(a+g\tan\theta\right)\le2\mu mg,
\]
故 $a\le g\left(2\mu-\tan\theta\right)$。

当 $\tan\theta\le\mu$ 时，$2\mu-\tan\theta\ge\tan\theta$，最大加速度 $a_{\max}=g\tan\theta$，
\[
  F_{\max}=4ma_{\max}=4mg\tan\theta,
\]
B 错误。

当 $\mu<\tan\theta\le2\mu$ 时，$2\mu-\tan\theta<\tan\theta$，故
\[
  a_{\max}=g\left(2\mu-\tan\theta\right),\qquad
  F_{\max}=4mg\left(2\mu-\tan\theta\right),
\]
C 正确。

D．推力向右时，整体加速度 $a$ 向右。对 $B$，杆的分力向右、大小为 $mg\tan\theta$，摩擦力方向可左可右，由
\[
  ma=mg\tan\theta\pm f,\qquad |f|\le2\mu mg
\]
得
\[
  g\left(\tan\theta-2\mu\right)\le a\le g\left(\tan\theta+2\mu\right).
\]
对 $A$，$N_{A}-mg\tan\theta=ma$，得 $N_{A}=m\left(a+g\tan\theta\right)>0$，$A$ 始终压紧左壁，不构成额外限制。故
\[
  4mg\left(\tan\theta-2\mu\right)\le F\le4mg\left(\tan\theta+2\mu\right),
\]
D 正确。
\end{answer}

\begin{answer}
\gksolhead{11}{实验题}
\gkanswer{答案 (3) $T=\dfrac{t_{0}}{10}$；(5) 线性的；(6) A；(7) 见解析}
(3) 由题图乙，$0\sim t_{0}$ 内有 10 个完整波形，即 $t_{0}=10T$，故
\[
  T=\frac{t_{0}}{10}.
\]

(5) 逐组计算 $\dfrac{T^{2}}{m}$：
\[
  \frac{0.059}{0.015}\approx3.93,\quad
  \frac{0.099}{0.025}\approx3.96,\quad
  \frac{0.138}{0.035}\approx3.94,\quad
  \frac{0.178}{0.045}\approx3.96,\quad
  \frac{0.217}{0.055}\approx3.95,
\]
比值近似为常数，故 $T^{2}$ 与 $m$ 成线性关系（$T^{2}\propto m$）。

(6) 由 $T^{2}\propto m$ 得 $T\propto\sqrt{m}$；又 $T$ 的单位为 \si{\second}，$m$ 的单位为 \si{\kilogram}，$k$ 的单位为 \si{\newton\per\metre}=\si{\kilogram\per\second\squared}。检验各表达式的单位：
\[
  \sqrt{\frac{m}{k}}:\ \sqrt{\frac{\si{\kilogram}}{\si{\kilogram\per\second\squared}}}
  =\si{\second},
\]
与周期的单位相符，A 正确；B、C、D 的单位或与 $T\propto\sqrt{m}$ 的关系不符，均错误。

(7) 可能产生误差的原因：弹簧并非轻质，其自身质量不可忽略；手机内部的磁体、导体使振动过程中产生电磁阻尼；钩码与小磁体未严格在同一竖直线上，振动不在竖直方向。
\end{answer}

\begin{answer}
\gksolhead{12}{实验题}
\gkanswer{答案 (1) 1000；(2) $R_{F}=\dfrac{R_{1}R_{3}}{R_{2}}$；(3) 见解析；(4) $1.75\times10^{-2}$；(5) $>$}
(1) 欧姆表用“$\times100$”挡，指针指在 \SI{10}{\ohm} 处，故
\[
  R_{F}=10\times100\,\si{\ohm}=1000\,\si{\ohm}.
\]

(2) 电压传感器示数为 0 时，$C$、$D$ 两点等势，桥路平衡。设流过 $R_{1}$、$R_{F}$ 的电流为 $I_{1}$，流过 $R_{2}$、$R_{3}$ 的电流为 $I_{2}$，则
\[
  I_{1}R_{1}=I_{2}R_{2},\qquad I_{1}R_{F}=I_{2}R_{3}.
\]
两式相除得
\[
  R_{F}=\frac{R_{1}R_{3}}{R_{2}}.
\]

(3) 根据表中数据描点，图线为一条过原点的直线，斜率约为 \SI{112}{\milli\volt\per\gram}。

\gkfig{figures/s12-graph}

(4) 由 $U$--$m$ 图线，$U=\SI{200}{\milli\volt}$ 时
\[
  m=\frac{200}{112}\,\si{\gram}\approx1.79\,\si{\gram}
   =1.79\times10^{-3}\,\si{\kilogram},
\]
故
\[
  F_{0}=mg=1.79\times10^{-3}\times9.8\,\si{\newton}
      \approx1.75\times10^{-2}\,\si{\newton}.
\]

(5) 非理想毫伏表内阻有限，接入 $C$、$D$ 间相当于在两端并联一个电阻，使 $C$、$D$ 间的电压小于同一压力下电压传感器（内阻无穷大）的示数。现读数仍为 \SI{200}{\milli\volt}，说明此时桥路的开路电压大于 \SI{200}{\milli\volt}，即传感器阻值偏离平衡值更多、所受压力更大，故 $F_{1}>F_{0}$。

\gkfig{figures/s12-loading}
\end{answer}

\begin{answer}
\gksolhead{13}{解答题}
\gkanswer{答案 (1) $p_{1}=\dfrac{V_{0}}{V_{0}+V_{1}}p_{0}$；(2) $\Delta F=\left[1-\left(\dfrac{V_{0}}{V_{0}+V_{1}}\right)^{n}\right]p_{0}S$}
(1) 第 1 次抽气时，助力气室中的气体等温膨胀，体积由 $V_{0}$ 变为 $V_{0}+V_{1}$。以原助力气室中封闭的气体为研究对象，由玻意耳定律
\[
  p_{0}V_{0}=p_{1}\left(V_{0}+V_{1}\right),
\]
解得
\[
  p_{1}=\frac{V_{0}}{V_{0}+V_{1}}p_{0}.
\]

(2) 每次抽气后，助力气室中气体的压强变为原来的 $\dfrac{V_{0}}{V_{0}+V_{1}}$，故第 $n$ 次抽气后
\[
  p_{n}=\left(\frac{V_{0}}{V_{0}+V_{1}}\right)^{n}p_{0}.
\]
对助力活塞与连杆 $AB$ 构成的系统，外部大气压力 $p_{0}S$ 向右，气室内气体压力 $p_{n}S$ 向左，驾驶员省力的大小
\[
  \Delta F=\left(p_{0}-p_{n}\right)S
  =\left[1-\left(\frac{V_{0}}{V_{0}+V_{1}}\right)^{n}\right]p_{0}S.
\]
\end{answer}

\begin{answer}
\gksolhead{14}{解答题}
\gkanswer{答案 (1) $v_{0}=\dfrac{2mgR\sin\theta}{B^{2}L^{2}}$；(2) $a_{0}=2g\sin\theta$}
\gkanswer{(3) $v=gt_{0}\sin\theta+\dfrac{mgR\sin\theta}{B^{2}L^{2}}$，\ $\Delta x=\dfrac{2m^{2}R^{2}g\sin\theta}{B^{4}L^{4}}$}
(1) 棒 a 与棒 b 构成闭合回路，b 静止，a 切割磁感线，相当于电源。设 a 匀速运动时的速度为 $v_{0}$，则
\[
  E=BLv_{0},\qquad I=\frac{E}{2R},
\]
a 所受安培力
\[
  F_{A}=BIL=\frac{B^{2}L^{2}v_{0}}{2R}.
\]
由平衡条件 $mg\sin\theta=F_{A}$ 得
\[
  v_{0}=\frac{2mgR\sin\theta}{B^{2}L^{2}}.
\]

(2) 棒 b 刚释放时，a 的速度仍为 $v_{0}$，回路中电流
\[
  I=\frac{BLv_{0}}{2R},\qquad BIL=\frac{B^{2}L^{2}v_{0}}{2R}=mg\sin\theta.
\]
b 所受安培力沿导轨向下。对 b 由牛顿第二定律
\[
  mg\sin\theta+BIL=ma_{0},
\]
解得
\[
  a_{0}=2g\sin\theta.
\]

(3) 棒 b 释放后，a、b 组成的系统所受外力沿导轨方向的总和为 $2mg\sin\theta$。经过时间 $t_{0}$ 两棒速度均为 $v$，由动量定理
\[
  2mg\sin\theta\cdot t_{0}=2mv-mv_{0},
\]
解得
\[
  v=gt_{0}\sin\theta+\frac{v_{0}}{2}
   =gt_{0}\sin\theta+\frac{mgR\sin\theta}{B^{2}L^{2}}.
\]
对 b 由动量定理
\[
  mg\sin\theta\cdot t_{0}+I_{F}=mv,
\]
其中安培力的冲量 $I_{F}=\sum BLi\Delta t=BLq$，而
\[
  q=\frac{\Delta\Phi}{2R}=\frac{BL\Delta x}{2R}.
\]
代入得
\[
  BL\cdot\frac{BL\Delta x}{2R}
  =m\left(gt_{0}\sin\theta+\frac{mgR\sin\theta}{B^{2}L^{2}}\right)-mg t_{0}\sin\theta
  =\frac{m^{2}gR\sin\theta}{B^{2}L^{2}},
\]
故
\[
  \Delta x=\frac{2m^{2}R^{2}g\sin\theta}{B^{4}L^{4}}.
\]
\end{answer}

\begin{answer}
\gksolhead{15}{解答题}
\gkanswer{答案 (1) $v_{1}=\sqrt{\dfrac{2m^{2}gb}{m_{0}(m_{0}+m)}}$，$x_{1}=\dfrac{ma}{m_{0}+m}$}
\gkanswer{(2) 见解析；\ (3) $v=\sqrt{\dfrac{4gb^{2}}{a+3b}}$}
(1) 小球与凹槽组成的系统在水平方向不受外力，水平方向动量守恒；系统机械能守恒。设小球第一次到最低点时凹槽的速度为 $v_{1}$（向右），小球的速度为 $v_{2}$（向左），则
\[
  m_{0}v_{1}=mv_{2},
\]
\[
  mgb=\frac{1}{2}m_{0}v_{1}^{2}+\frac{1}{2}mv_{2}^{2}.
\]
解得
\[
  v_{1}=\sqrt{\frac{2m^{2}gb}{m_{0}\left(m_{0}+m\right)}}.
\]
设该过程中凹槽向右移动 $x_{1}$，小球相对地面向左移动 $x_{2}$，则
\[
  m_{0}x_{1}=mx_{2},\qquad x_{1}+x_{2}=a,
\]
解得
\[
  x_{1}=\frac{ma}{m_{0}+m}.
\]

(2) 设某时刻小球位于 $(x,\,y)$，凹槽中心位于 $x_{1}'$，小球相对凹槽中心的水平坐标为 $x-x_{1}'$。初始时系统的水平总动量为零，故任意时刻
\[
  m_{0}x_{1}'+mx=ma,
\]
即
\[
  x_{1}'=\frac{m}{m_{0}}\left(a-x\right).
\]
小球在凹槽内的半椭圆轨道上，满足
\[
  \frac{\left(x-x_{1}'\right)^{2}}{a^{2}}+\frac{y^{2}}{b^{2}}=1\qquad(y\le0).
\]
代入 $x_{1}'$ 得
\[
  x-x_{1}'=x-\frac{m}{m_{0}}\left(a-x\right)
  =\frac{\left(m_{0}+m\right)x-ma}{m_{0}},
\]
故小球的轨迹方程为
\[
  \frac{\left(x-\frac{m}{m_{0}+m}a\right)^{2}}
       {\left(\frac{am_{0}}{m_{0}+m}\right)^{2}}
  +\frac{y^{2}}{b^{2}}=1\qquad(y\le0).
\]

\gkfig{figures/s15-coords}

(3) 当 $\dfrac{m_{0}}{m}=\dfrac{b}{a-b}$ 时，轨迹方程化为
\[
  \frac{\left[x-(a-b)\right]^{2}}{b^{2}}+\frac{y^{2}}{b^{2}}=1,
\]
即小球的轨迹是圆心为 $(a-b,\,0)$、半径为 $b$ 的圆（下半圆）。

\gkfig{figures/s15-circle}

小球下降 $h=\dfrac{b}{2}$ 时，由几何关系，小球与圆心的连线与水平方向夹角 $\alpha$ 满足
\[
  \sin\alpha=\frac{b/2}{b}=\frac{1}{2},\qquad \alpha=30^{\circ}.
\]
速度沿圆的切线方向，与水平方向成 $60^{\circ}$。

水平方向动量守恒：
\[
  mv\cos60^{\circ}=m_{0}v'.
\]
机械能守恒：
\[
  mg\cdot\frac{b}{2}=\frac{1}{2}mv^{2}+\frac{1}{2}m_{0}v'^{2}.
\]
由 $mv\cos60^{\circ}=m_{0}v'$ 得 $v'=\dfrac{m}{2m_{0}}v$，代入能量式：
\[
  \frac{gb}{2}=\frac{1}{2}v^{2}+\frac{m^{2}}{8m_{0}}v^{2}
  =\frac{1}{2}v^{2}\left(1+\frac{m}{4m_{0}}\right).
\]
由 $\dfrac{m}{m_{0}}=\dfrac{a-b}{b}$ 得
\[
  1+\frac{m}{4m_{0}}=1+\frac{a-b}{4b}=\frac{a+3b}{4b},
\]
故
\[
  \frac{gb}{2}=\frac{1}{2}v^{2}\cdot\frac{a+3b}{4b},
\]
解得
\[
  v=\sqrt{\frac{4gb^{2}}{a+3b}}.
\]
\end{answer}
